\paragraph{Conjugate-coordinate class: the Zeeman-dressed vector sector at order \(\epsilon^2\), its logarithmic coefficient, the spin blocks and the continuation (conditional; bounded no-go; correction of the sign of a logarithmic term in the preceding finite-part values).}
Keep the setting of the preceding paragraphs: Poisson limit, leaf \(\mathbb{CP}^{1,2}\), the nilpotent (leading) symbol with centre flux \([\omega_a,\omega_b]=\mathrm iF_{ab}\), the central plane-wave field with \(P=y^c\omega_c+c\) at the seam, \(A^a=\epsilon\alpha^a\cos\phi\), \(\kappa_a=W_a(\phi)\), and the species data of the paragraph on the supertrace at an arbitrary matrix fluctuation. With \(\Omega_a=\omega_a-\epsilon\alpha_aP\sin\phi\) one has \([\Omega_a,\Omega_b]=\mathrm iF_{ab}+\mathrm i\epsilon P\cos\phi\,K_{ab}\), \(K_{ab}=\kappa_a\alpha_b-\kappa_b\alpha_a\), and the vector operator is \(\eta^{ab}\Omega_a\Omega_b\otimes1+2[\Omega,\Omega]\eta+\mu^2\) (the coefficient \(-2\mathcal Z_{ab}\) is the one for which the longitudinal modes \(D_a\lambda\) are mapped to \(D_a\) of the scalar operator, with \(\mathcal Z_{ab}=-[\Omega_a,\Omega_b]\)). At \(\epsilon=0\) the spin factor \(\mathrm e^{-\mathrm is Z}=\mathrm e^{2sF\eta}\) has the four blocks \(\mathrm e^{\pm2sf_e}\) and \(\mathrm e^{\pm2\mathrm isf_m}\) (trace \(\chi_V\) of the preceding paragraphs). The order \(\epsilon^2\) of \(\operatorname{Tr}\mathrm e^{-\mathrm is(H+Z)}\) was computed by the double Duhamel integral of the preceding paragraph with three kinds of insertion: the kinetic one, the Zeeman one \(\epsilon P\cos\phi\,N\), \(N=2\mathrm iK\eta\) for the vector and \(\tfrac{\mathrm i}{4}[\Gamma^a,\Gamma^b]K_{ab}\) for the Dirac block, and their products, with the spin factors \(\operatorname{tr}[\mathrm e^{2sF\eta}N]\) and \(\operatorname{tr}[\mathrm e^{(s-\tau)K_F}N\mathrm e^{\tau K_F}N]\), \(K_F=2F\eta\). For a scalar species the result equals the preceding engine (\(10^{-13}\)); the response of every species is invariant under \(\alpha\to\alpha+\lambda\kappa\) (\(10^{-14}\) at \(s=0.4\), \(10^{-8}\) at complex \(s\), limited by quadrature).

\emph{Logarithmic coefficient.} Write \(16\pi^2c=\lim_{s\to0}16\pi^2\,s\,T_2(s)\), \(F_1=y^2\bigl((\kappa\cdot\alpha)^2-\kappa^2\alpha^2\bigr)\), \(F_2=\kappa^2(\alpha\cdot y)^2\), so that the form \(A\) of the next-to-leading-symbol paragraph is \(A=(F_1-2F_2)/24\). Fitting the ten independent components of the quadratic form at one orbit point gives, for the four-dimensional vector with the Zeeman dressing, \(16\pi^2c_V=4A-F_1\) (coefficients \(g_1=-0.8337\), \(g_2=-0.3333\) against \(-5/6\), \(-1/3\)), and for one \(4\times4\) Dirac block \(16\pi^2c_D=4A-\tfrac12F_1\) (\(-0.33346\), \(-0.33333\) against \(-1/3\), \(-1/3\)); the scalar value is \(A\). The extra term \(-F_1\) (\(-\tfrac12F_1\)) is the term quadratic in \(N\), proportional to \(\operatorname{tr}\Sigma\Sigma\), with \(\operatorname{tr}\Sigma_D\Sigma_D=\tfrac12\operatorname{tr}\Sigma_V\Sigma_V\). The weighted sum over the supertrace content (\(\sum_iw_ic_i\) with the weights of the preceding paragraph) vanishes to \(10^{-5}\) at the three sampled points, so the massless resummed character term has no ultraviolet logarithm at this order. The mass-linear term has the logarithmic form \(M_{\mathrm{str}}A-\mu^2F_1\): for \(\mu^2\ne0\) it is not proportional to \(A\), and since both \(F_1\) and \(F_2\) contain \(y^\mu\) through \(y^2\) and \(\alpha\cdot y\), the dependence on \(y^\mu\) shown in the next-to-leading-symbol paragraph is not removed by the Zeeman term.

\emph{Spin blocks and the continuation.} Decomposing the vector response in the eigenbasis \(E_\pm\), \(B_\pm\) of \(K_F\) (diagonal blocks and pair blocks with the interpolating phase \(\mathrm e^{(s-\tau)\mu_k+\tau\mu_l}\)), at one orbit point the diagonal and pair blocks containing \(E_+\) grow along the rays \(s=|s|\mathrm e^{-\mathrm i\vartheta}\) with \(\vartheta<\vartheta_*=\arctan(f_e/f_m)\) and those containing \(B_+\) grow for \(\vartheta>\vartheta_*\); at \(\vartheta_*\) every block grows polynomially. No ray gives a convergent proper-time integral of the Zeeman-dressed term. A continuation is therefore a hypothesis. Two prescriptions were compared on the leading-symbol integral \(\int\mathrm ds\,f_ef_m[\sinh(f_es)\sin(f_ms)]^{-1}(\chi_V-4)\) at \((f_e,f_m)=(1.6,0.9)\): the groups \(\{E_\pm\}\) and \(\{B_\pm\}\) integrated on their own rays (\(-4.007-9.981\mathrm i\) and \(-4.015-9.973\mathrm i\) for two pairs of rays), and the integral on \(\vartheta_*\) with the factor \(\mathrm e^{-\mathrm ism^2}\), \(m^2>0\), extrapolated to \(m^2\to0^+\) (\(-4.007\) to \(-4.013\), \(-9.984\) to \(-9.992\mathrm i\)); they agree to the stated accuracy. The imaginary part is the unstable branch (a Nielsen--Olesen-type mode) and is not removed.

\emph{Finite part at one orbit point.} For the tetrad part (\(c=0\)) at one orbit point and one centre direction (flux rescaled to \(f_e=1.004\), \(f_m=0.396\)), the finite part \(\mathrm{Fin}=\lim[\int_{s_0}^\infty T_2\,\mathrm ds-c\ln(1/s_0)]\) of the vector term, the scalar term and the massless supertrace were evaluated on \(\vartheta_*\) with the factor \(\mathrm e^{-\mathrm ism^2}\) for \(m^2=0.4,\dots,1.3\) in units of \(\max(f_e,f_m)\) (quadrature checked against the number of nodes; tail truncation \(<10^{-4}\)); the extrapolation \(m^2\to0\) is not stable (entries change by order one between orders), so no value at \(m^2=0\) is claimed. In this regulated family: (i) all three forms are gauge invariant (\(|Q\kappa|/|Q||\kappa|\le6\cdot10^{-6}\)); (ii) the vector form is not proportional to its logarithmic form \(4A-F_1\) (relative defect \(0.91\)--\(0.98\)) nor to \(A\) (\(0.79\)--\(0.82\)), and the massless supertrace form is not proportional to \(A\) (\(0.87\)--\(1.0\)), a statement independent of the logarithm scheme for the supertrace, which has no logarithm; (iii) on the quotient by \(\kappa\) the real part of the vector form is positive definite for every \(m^2\) tried, its imaginary part is positive definite for \(m^2\le0.6\) and indefinite for \(m^2\ge0.8\), the scalar form is indefinite in the real part, and the real part of \(x\,Q_S-3Q_V\), \(x=n_{\mathrm{ex}}\mu_i^2\), is definite on the quotient for every \(x\) in an interval containing \([0,21]\) at every \(m^2\) tried; (iv) the sign statements for the vector and the massless forms change with \(m^2\) and with the real or imaginary part, and a sign statement for a quotient spectrum is not invariant under the finite shift \(c\ln\lambda\) of the scheme for forms that carry a logarithm. The result is for one orbit point.

\emph{Correction to the preceding paragraph.} The routine that produced the finite-part values of the preceding paragraph added the term \(-c\ln s_1\) where the definition of \(\mathrm{Fin}\) gives \(+c\ln s_1\) (\(s_1\) the matching point; a test function \(c\,\mathrm e^{-s}/s\) returns \(5.41+1.20\mathrm i\) instead of \(-\gamma\)). The values differ by \(2c\ln s_1\), a multiple of the form \(A\). The non-proportionality to \(A\) is therefore unchanged: the defect \(\|Q-\lambda A\|/\|A\|\), which is shift invariant, is \(0.19\) and \(0.075\) (real, imaginary parts; first sample), \(4.5\cdot10^{-3}\) and \(4.4\cdot10^{-3}\) (second sample) and \(2.8\cdot10^{-4}\) and \(8\cdot10^{-8}\) (strongly electric sample, \(f_e/f_m\simeq42\)), whereas the normalised residuals quoted there become \(0.63\) and \(0.82\), \(0.49\) and \(0.99\), \(0.011\) and \(0.94\); the quotient spectra of the real parts remain of both signs.

This is a statement for the Poisson limit, one loop, the nilpotent (leading) symbol resummed in the flux with the central seam field, the Lorentzian \(\mathrm i\epsilon\) proper time, the tetrad part at one orbit point and one centre direction, a regulator \(\mathrm e^{-\mathrm ism^2}\) not removed, the signs of the Zeeman couplings fixed by the longitudinal-mode identity above, and numerical evaluation. It does not include the charge and cross parts, the window average over the centre variable, the weight-one terms, the continuation \(m^2\to0\) of the four-dimensional integral, the sign convention of the factor \(\mathrm i\) of the mass-linear term, or a derivation of the continuation prescription, and it identifies no induced form with the emergent metric of Chapter~\ref{ch:vierbein-emergence}, makes no Wick-rotation claim, and gives no value or sign of \(G_{\mathrm{eff}}\). This changes no ledger label; Problem~F remains open.

